\documentclass[10pt,a4paper]{amsart}
\usepackage[margin=19mm]{geometry}
\usepackage{amsmath,amssymb,mathtools}
\usepackage[scr=rsfs,cal=euler]{mathalfa}
\usepackage{xfrac}
\usepackage[unicode,hidelinks,bookmarks=false]{hyperref}
\newcommand{\ie}{i.e.\ }
\newcommand{\cf}{cf.\ }
\newcommand{\resp}{respectively\ }
\newcommand{\Subsection}{\subsection}
\providecommand{\nobreakdash}{\nobreak\hbox{-}}
\setlength{\emergencystretch}{2em}
\multlinegap0pt
\usepackage{amssymb}
\newtheorem{lemma}{Lemma}
\newcommand{\Z}{\mathbb{Z}}
\newcommand{\abs}[1]{\lvert#1\rvert}
\title{A Fourier-Free Density-Increment Proof of Roth's Theorem}
\author{Mark Lewko}
\date{Source: arXiv:2605.19310v1 (CC BY 4.0)}
\begin{document}
\maketitle

\noindent\emph{Source and licence.} A Fourier-Free Density-Increment Proof of Roth's Theorem. Version arXiv:2605.19310v1 (CC BY 4.0), distributed by arXiv under the Creative Commons Attribution 4.0 International licence, http://creativecommons.org/licenses/by/4.0/, as recorded in the arXiv licence metadata for this exact version and shown on the abstract page https://arxiv.org/abs/2605.19310v1. Full licence notice: \texttt{licenses/arxiv-2605.19310-CC-BY-4.0.txt}.\par\smallskip

\noindent\textit{Editorial scope.} Counted unit: the article's lemma lem:lambda-energy (source0.tex lines 146--153). The article's complete original proof of it is delivered with it (source0.tex lines 155--199, one \texttt{proof} environment of 45 lines). No mathematics is written, repaired or re-derived: the statement and the proof are the article's own lines, in the article's own notation, and no retained line is substituted or re-flowed.  No further source-local context is needed: every cross-reference of every retained line resolves inside the two delivered spans.  Nothing was omitted from any retained span, so the package carries no omission record.  No retained line contains a citation, so the wrapper carries no bibliography and no bibliography entry.  No retained line contains a figure environment, an image-inclusion command or drawing code, so no image is delivered and the package declares no asset dependency.  Editorial formatting only, in addition: an AMS \texttt{amsart} preamble that restates the article's own declarations and macros used by the retained lines, namely the environments \texttt{lemma} on separate counters, together with the standard packages those lines need.  The retained environments print in this document as Lemma 1 in order of appearance, whereas the article prints the same items with the numbering of its own full document.  The complete native source of the article as distributed in the arXiv e-print for this exact version is bundled unchanged as \texttt{sources/200909/source0.tex}.
\begin{lemma}\label{lem:lambda-energy}
Let $g_1,g_2,g_3\colon\Z_m\to\mathbb R$ be bounded in absolute value by an
absolute constant. Then
\[
        \abs{\Lambda(g_1,g_2,g_3)}^4
        \lesssim m^5\min_{1\leq i\leq3}\mathcal E(g_i).
\]
\end{lemma}

\begin{proof}
It suffices to prove the following estimate. Let $F,h_1,h_2\colon\Z_m\to\mathbb R$,
with $h_1,h_2$ bounded, and let $a,b\in\Z_m^\times$ with $a\neq b$. Then
\[
        \left|\sum_{y,r\in\Z_m}F(y)h_1(y+ar)h_2(y+br)\right|^4
        \lesssim m^5\mathcal E(F).
\]
The three possible choices of the distinguished slot in $\Lambda$ are obtained,
after a change of variables, from this estimate by taking $(a,b)=(1,2)$,
$(-1,1)$, and $(-2,-1)$.

Let $T$ denote the sum inside the fourth-power. Put $\lambda=ba^{-1}$ and
change variables from $r$ to $u=y+ar$. Then
\[
T=\sum_{y,u\in\Z_m}F(y)h_1(u)h_2((1-\lambda)y+\lambda u).
\]
By Cauchy--Schwarz in $u$,
\[
\abs{T}^2\lesssim m\sum_{u\in\Z_m}
\left|\sum_{y\in\Z_m}F(y)h_2((1-\lambda)y+\lambda u)\right|^2.
\]
Expanding the square, then putting $t=y'-y$ and $z=(1-\lambda)y+\lambda u$, gives
\[
\begin{aligned}
\abs{T}^2
&\lesssim m\sum_{t\in\Z_m}
\left(\sum_{y\in\Z_m}F(y)F(y+t)\right) \\
&\qquad\qquad\times
\left(\sum_{z\in\Z_m}h_2(z)h_2(z+(1-\lambda)t)\right).
\end{aligned}
\]
A final Cauchy--Schwarz in $t$ gives
\[
\begin{aligned}
\abs{T}^2
&\lesssim m\mathcal E(F)^{1/2}
\left(\sum_{t\in\Z_m}
\left(\sum_{z\in\Z_m}h_2(z)h_2(z+(1-\lambda)t)\right)^2
\right)^{1/2} \\
&\lesssim m^{5/2}\mathcal E(F)^{1/2},
\end{aligned}
\]
using that $|h_2|\leq 1$. Squaring proves the estimate. Applying this with each
$g_i$ as the distinguished function gives the minimum on the right.
\end{proof}

\end{document}
